% !TEX root = ../../main_docker.tex
% ---------------------------------------------------------------------
%  Navigation: "Wenn die Aufgabe lautet ..." -- von der Aufgabenstellung
%  direkt zur Abschnittsnummer. Steht vor dem Inhaltsverzeichnis.
% ---------------------------------------------------------------------

\begin{tcolorbox}[zfbox, colback=black!2, colframe=black!45,
  coltitle=white, colbacktitle=black!45, title={Wenn die Aufgabe lautet \dots}]
\begin{tabularx}{\linewidth}{@{}Xr@{}}
Kommando zu einer Situation nennen & \ref{sec:ref} \\
Image- und Container-Grösse rechnen & \ref{rez:groesse} \\
Was passiert beim Beenden/Trennen & \ref{sec:beenden} \\
Wohin schreibt der Container & \ref{sec:overlay} \\
Volume-Art zu einer Anforderung & \ref{rez:volume} \\
Was ein Tag-Name ausgeschrieben heisst & \ref{rez:tag} \\
Dockerfile zu einer Aufgabe schreiben & \ref{sec:dockerfile} \\
Anzahl Layer eines Dockerfiles & \ref{rez:layer} \\
Was \R{docker run} tatsächlich startet & \ref{form:entrypoint} \\
Warum ein Image so gross/klein ist & \ref{sec:multistage} \\
Mehrstufigen Build schreiben & \ref{rez:multistage} \\
Statisch oder dynamisch binden & \ref{sec:linken} \\
\end{tabularx}
\vspace{2pt}
\textcolor{accent}{\rule{7pt}{7pt}}\,Formel \;
\textcolor{red!65}{\rule{7pt}{7pt}}\,Falle \;
\textcolor{green!40!black}{\rule{7pt}{7pt}}\,Rezept \;
\textcolor{black!55}{\rule{7pt}{7pt}}\,Beispiel \;
-- der letzte Abschnitt ist das Stichwortregister.
\end{tcolorbox}
